% ======================================================================
% Ontology of Continua — Core 1.3
% Cross-K module: crossk_k1_k2.tex
% Cross-level relations between K1 and K2
% ======================================================================

\subsubsection{Межуровневая структура K1–K2}
\label{sec:crossk-k1-k2}

Переход $K_1 \rightarrow K_2$ обозначает возникновение геометрической и
энергетической структуры из минимального представительного континуума.
Если $K_1$ содержит одну ось, примитивные потенциалы и тривиальные потоки,
то $K_2$ вводит:
\begin{itemize}
  \item несколько осей, формирующих координатную структуру;
  \item явные соотношения потенциал–градиент;
  \item нетривиальные потоки $J_2$;
  \item причинно-следственный порядок;
  \item более богатую область $\Omega(K_2)$, поддерживающую физическую
        динамику.
\end{itemize}

Межуровневая структура фиксирует, как простая представительная геометрия
$K_1$ расширяется в физико-подобную геометрию $K_2$.

% ----------------------------------------------------------------------
\subsubsection{1. Задействованные уровни и их роли}

\paragraph{K1.}
$K_1$ предоставляет:
\begin{itemize}
  \item одну ось $A_1$;
  \item примитивный потенциал $P_1$;
  \item базовую структуру границы $\partial\Omega(K_1)$;
  \item первые потоки на основе градиента $J_1$;
  \item тривиальные циклы, но без полной геометрической или причинной
        структуры.
\end{itemize}

\paragraph{K2.}
$K_2$ расширяет размерность и вводит:
\begin{itemize}
  \item хотя бы одну новую ось $A_2$;
  \item многомерную координатную карту;
  \item энергетические потенциалы и поля;
  \item нетривиальные динамические потоки $J_2$;
  \item причинный порядок и физически интерпретируемые ограничения;
  \item более богатое пространство состояний $\Omega(K_2)$ с нетривиальной
        топологией.
\end{itemize}

K₂ — первый континуум, поддерживающий перколяцию, распространение полей и
динамическое поведение.

% ----------------------------------------------------------------------
\subsubsection{2. Совместные и наследуемые оси и пороги}

\paragraph{Оси.}
Ось $A_1$ уровня $K_1$ становится одной из координатных осей внутри
$A(K_2)$:
\[
A(K_1) \subset A(K_2) = \{A_1, A_2, \dots\}.
\]
Новая ось $A_2$ позволяет возникнуть градиентам, полям и переносу через
многомерную область.

\paragraph{Пороги.}
Порог существования $\Theta_1$ обеспечивает устойчивость $A_1$; при его
нарушении весь переход разрушается:
\[
\Theta_1^{\mathrm{death}} \Rightarrow \Omega(K_2)=\varnothing.
\]

В $K_2$ вводятся новые классы порогов:
\begin{itemize}
  \item \textbf{энергетические пороги} для полей и потенциалов,
  \item \textbf{пороги градиента} для ограничения потоков,
  \item \textbf{пороги перколяции} для связности,
  \item \textbf{пороги причинности} для поддержания порядка событий.
\end{itemize}

Эти пороги уточняют и расширяют структуру, унаследованную от $K_1$.

% ----------------------------------------------------------------------
\subsubsection{3. Межуровневые потоки, циклы и напряжения}

\paragraph{Потоки.}
Потоки $J_2$ наследуются от $J_1$, но усложняются из-за дополнительной
оси и энергетических потенциалов:
\[
J_2 = \nabla P_2 + \text{transport terms along } A_2.
\]
Сюда входят диффузионные, полевые и причинные режимы распространения.

\paragraph{Циклы.}
Циклы в $K_2$ соответствуют замкнутым геометрическим или энергетическим
петлям:
\[
C_2 : \oint_{\gamma} J_2 \cdot dA.
\]
Они невозможны в $K_1$, где размерности недостаточно.
Поэтому $C_2$ является строгим вверх направленным расширением тривиального
цикла $K_1$.

\paragraph{Напряжение.}
Межуровневое напряжение определяется как:
\[
T_{1,2} = h(P_2 - P_1,\ \Theta_2 - \Theta_1,\ A_2).
\]
Если $T_{1,2}$ превышает размерностный порог, новая ось $A_2$ становится
неустойчивой, а $\Omega(K_2)$ возвращается к дегенеративной форме.

% ----------------------------------------------------------------------
\subsubsection{4. Условия рождения и исчезновения между уровнями}

\paragraph{Рождение \texorpdfstring{$K_2$}{K_2}.}
$K_2$ возникает, когда:
\[
T_1 > \Theta_{1,\mathrm{dim}}
\quad \text{и} \quad
A_2 \in M_2 \setminus A(K_1),
\]
с соответствующим расширением пространства состояний:
\[
\Omega(K_2) = \Omega(K_1) \cup \Delta\Omega_{\mathrm{geom}}.
\]

\paragraph{Распространение гибели.}
Если $\Theta_1$ не выполняется, восходящего расширения не существует.
Если $\Theta_2$ не выполняется, существует только падение $K_2$; $K_1$
остаётся корректным, поскольку его структура не зависит от потенциалов и потоков
высшего уровня.

% ----------------------------------------------------------------------
\subsubsection{5. Концептуальные примеры}

\paragraph{Возникновение геометрии.}
Единственная ось ($K_1$) не поддерживает потоки с ротором, дивергенцией или
многонаправленным распространением. Добавление $A_2$ делает возможным:
\[
\mathrm{grad},\ \mathrm{div},\ \mathrm{curl}-\text{подобные структуры}.
\]

\paragraph{Перколяция.}
Многомерная область поддерживает переходы связности (порог перколяции $p_c$),
которые становятся частью пороговой структуры $K_2$, но не имеют аналога в
$K_1$.

\paragraph{Причинность.}
Динамическое распространение в разных направлениях в многомерной структуре
приводит к первому появлению причинного порядка: невозможному в $K_1$.

Эти примеры иллюстрируют существенный смысл перехода K1→K2:
\[
\textbf{Это рождение геометрии, полей и динамики из минимального
представительного континуума.}
\]
